% =====================================================================
\chapter{Soluciones de los ejercicios}
\label{ap:soluciones}
% =====================================================================

Cada solución se presenta en cinco partes: datos, planteamiento, desarrollo, resultado y comprobación. En los pasos intermedios se redondea a cuatro cifras significativas salvo que se indique otra cosa; los resultados exactos se dan como fracción cuando existe. Las pérdidas se expresan en nats.

% ---------------------------------------------------------------------
\section{Tema 1}

\begin{solucion}{eje:1-tipos}
\datos cuatro problemas; el modelo debe devolver probabilidades en los de clasificación.
\planteamiento la salida es real (regresión) o una de un conjunto finito de clases (\cref{def:supervisado}); para predicciones probabilísticas se usa la pérdida logarítmica (\cref{prop:logloss}).
\pasos
\begin{enumerate}[label=(\alph*)]
  \item Precio: número real, \textbf{regresión}; pérdida cuadrática.
  \item Dígito: \textbf{clasificación} con $K=10$ clases; pérdida logarítmica (entropía cruzada).
  \item Token siguiente: \textbf{clasificación} con $K=128\,256$ clases; pérdida logarítmica, la del preentrenamiento.
  \item Correo deseado o no: \textbf{clasificación binaria}, $K=2$; pérdida logarítmica, que para dos clases es la entropía cruzada binaria (\cref{tema:3}).
\end{enumerate}
\resultado \textbf{(a) regresión; (b), (c) y (d) clasificación con 10, $128\,256$ y 2 clases.}
\comprobacion la pérdida 0-1 solo sería adecuada si el modelo devolviese una única clase sin probabilidad.
\end{solucion}

\begin{solucion}{eje:1-particion}
\datos $1797$ ejemplos; 20\,\% de prueba, 20\,\% de validación y el resto de entrenamiento.
\planteamiento redondear $0{,}2\cdot 1797$ y restar.
\pasos
\begin{enumerate}[label=\arabic*.]
  \item $0{,}2\cdot 1797=359{,}4$, que se redondea a $359$ ejemplos de prueba y $359$ de validación.
  \item Entrenamiento: $1797-2\cdot 359=1797-718=1079$.
  \item Si $\lambda$ se elige maximizando la exactitud en prueba, la prueba interviene en el ajuste y, por la \cref{prop:optimismo}, su exactitud deja de ser una estimación insesgada del rendimiento: resulta optimista.
\end{enumerate}
\resultado \textbf{$1079$ de entrenamiento, $359$ de validación y $359$ de prueba.} El hiperparámetro se elige con validación; la prueba se usa una sola vez al final.
\comprobacion $1079+359+359=1797$.
\end{solucion}

\begin{solucion}{eje:1-bayes}
\datos el \cref{tab:conjunta}; $P(x=0)=0{,}4$, $P(x=1)=0{,}6$, $\E[y\mid x=0]=0{,}75$, $\E[y\mid x=1]=4/3$.
\planteamiento $\Var(y\mid x)=\E[y^{2}\mid x]-\E[y\mid x]^{2}$; el mejor constante es $c=\E[y]$ y su riesgo es $\Var(y)$ (\cref{teo:regresion} sin $x$).
\pasos
\begin{enumerate}[label=\arabic*.]
  \item $\E[y^{2}\mid x=0]=(0\cdot 0{,}2+1\cdot 0{,}1+4\cdot 0{,}1)/0{,}4=0{,}5/0{,}4=1{,}25$; $\Var(y\mid x=0)=1{,}25-0{,}75^{2}=1{,}25-0{,}5625=0{,}6875$.
  \item $\E[y^{2}\mid x=1]=(0{,}2+4\cdot 0{,}3)/0{,}6=1{,}4/0{,}6=7/3$; $\Var(y\mid x=1)=7/3-16/9=5/9\approx 0{,}5556$.
  \item (a) $\E[\Var(y\mid x)]=0{,}4\cdot 0{,}6875+0{,}6\cdot 5/9=0{,}275+0{,}3333=0{,}6083$.
  \item (b) $\E[y]=0{,}1+0{,}2+2(0{,}1+0{,}3)=1{,}1$; $\E[y^{2}]=0{,}1+0{,}2+4(0{,}1+0{,}3)=1{,}9$; $\Var(y)=1{,}9-1{,}21=0{,}69$.
  \item (c) $\E[(f^{*}(x)-1{,}1)^{2}]=0{,}4(0{,}75-1{,}1)^{2}+0{,}6(4/3-1{,}1)^{2}=0{,}4\cdot 0{,}1225+0{,}6\cdot 0{,}05444=0{,}049+0{,}03267=0{,}08167$.
\end{enumerate}
\resultado \textbf{(a) $0{,}6083$; (b) $0{,}69$; (c) la diferencia es $0{,}0817$.}
\comprobacion $0{,}6083+0{,}0817=0{,}6900$, como predice el \cref{teo:regresion} con $f=1{,}1$. Además, $\E[f^{*}(x)]=0{,}4\cdot 0{,}75+0{,}6\cdot 4/3=0{,}3+0{,}8=1{,}1=\E[y]$.
\end{solucion}

\begin{solucion}{eje:1-constante}
\datos $x=(1,2,4,9,14)$, $n=5$.
\planteamiento derivar cada suma respecto a $b$; la derivada de $\abs{x_i-b}$ es $\sign(b-x_i)$ si $b\ne x_i$.
\pasos
\begin{enumerate}[label=\arabic*.]
  \item (a) $\frac{d}{db}\sum(x_i-b)^{2}=-2\sum(x_i-b)=0$, luego $b=\frac1n\sum x_i=30/5=6$, la media. La segunda derivada, $2n>0$, confirma que es un mínimo.
  \item (b) $\frac{d}{db}\sum\abs{x_i-b}=\#\{x_i<b\}-\#\{x_i>b\}$. Es negativa mientras haya más datos a la derecha que a la izquierda de $b$ y positiva en el caso contrario; cambia de signo en $b=4$, la mediana.
  \item (c) $\sum(x_i-6)^{2}=25+16+4+9+64=118$ y $\sum(x_i-4)^{2}=9+4+0+25+100=138$; $\sum\abs{x_i-4}=3+2+0+5+10=20$ y $\sum\abs{x_i-6}=5+4+2+3+8=22$.
  \item Si 14 se sustituye por un valor mayor, la media crece con él, pero la mediana sigue siendo~4.
\end{enumerate}
\resultado \textbf{la media, $b=6$, para la pérdida cuadrática; la mediana, $b=4$, para la absoluta}, que es menos sensible al valor atípico.
\comprobacion cada constante gana en su propia pérdida: $118<138$ y $20<22$.
\end{solucion}

\begin{solucion}{eje:1-logloss}
\datos $p$ con $p_k>0$ y $\sum_kp_k=1$; $q$ en el símplex.
\planteamiento (a) lagrangiana con la restricción $\sum_kq_k=1$; (b) desarrollar el cuadrado.
\pasos
\begin{enumerate}[label=\arabic*.]
  \item (a) $\Lambda(q,\mu)=-\sum_kp_k\ln q_k+\mu(\sum_kq_k-1)$. Entonces $\partial\Lambda/\partial q_k=-p_k/q_k+\mu=0$, luego $q_k=p_k/\mu$; la restricción da $\mu=1$ y $q=p$.
  \item Es el mínimo global: la función es estrictamente convexa en el símplex (sus segundas derivadas son $p_k/q_k^{2}>0$ y las cruzadas, nulas) y tiende a $+\infty$ si algún $q_k\to 0$. Equivalentemente, por la \cref{prop:logloss}, el exceso sobre el mínimo es $D_{\mathrm{KL}}(p\Vert q)\ge 0$.
  \item (b) $\norm{q-e_k}^{2}=\norm{q}^{2}-2q_k+1$. Ponderando con $p_k$ y sumando: $\norm{q}^{2}-2q^{\top}p+1=\norm{q-p}^{2}+1-\norm{p}^{2}$. El único término que depende de $q$ es $\norm{q-p}^{2}$, mínimo en $q=p$.
  \item (c) Logarítmica: $H(p)=0{,}5\ln 2+2\cdot 0{,}25\ln 4=0{,}5\cdot 0{,}6931+0{,}5\cdot 1{,}3863=1{,}040$ nats.
  \item Brier: $1-\norm{p}^{2}=1-(0{,}25+0{,}0625+0{,}0625)=0{,}625$.
\end{enumerate}
\resultado \textbf{ambos riesgos se minimizan en $q=p$; los mínimos son $H(p)=1{,}040$ nats y $0{,}625$.} $\blacksquare$
\comprobacion con $q=(1/3,1/3,1/3)$ el riesgo logarítmico es $\ln 3=1{,}099>1{,}040$ y el de Brier, $\norm{q-p}^{2}+0{,}625=0{,}0417+0{,}625=0{,}667>0{,}625$.
\end{solucion}

% ---------------------------------------------------------------------
\section{Tema 2}

\begin{solucion}{eje:2-rango}
\datos $x=(0,1,2)$, $y=(0,1,3)$; columnas de $X$: $\mathbf{1}$, $x$, $x$.
\planteamiento \cref{lem:rango} y \cref{cor:unicidad}.
\pasos
\begin{enumerate}[label=\arabic*.]
  \item (a) $X^{\top}X=\begin{pmatrix}3&3&3\\3&5&5\\3&5&5\end{pmatrix}$, con dos filas iguales: determinante nulo, \textbf{no es invertible}. Las columnas de $X$ son dependientes: $X(0,1,-1)^{\top}=0$.
  \item (b) Las predicciones son $b+(w_1+w_2)x$, y solo importa la suma $w_1+w_2$. Por el \cref{cor:unicidad}, los minimizadores son $\theta=(-1/6,\;w_1,\;3/2-w_1)$ con $w_1\in\R$ arbitrario: una recta de soluciones, todas con las mismas predicciones.
  \item (c) Primera solución: eliminar la columna repetida (o, en general, las columnas dependientes). Segunda: regularizar con ridge (\cref{tema:6}), que hace invertible $X^{\top}X+\lambda I$. Tercera: la pseudoinversa.
  \item Norma mínima: minimizar $w_1^{2}+(3/2-w_1)^{2}$ da $2w_1-2(3/2-w_1)=0$, es decir, $w_1=w_2=3/4$.
\end{enumerate}
\resultado \textbf{$X^{\top}X$ es singular; hay infinitos minimizadores $(-1/6,\,w_1,\,3/2-w_1)$, y el de norma mínima es $(-1/6;\;0{,}75;\;0{,}75)$.}
\comprobacion el cálculo numérico de $X^{+}y$ da $(-0{,}1667;\,0{,}75;\,0{,}75)$, y $0{,}75+0{,}75=1{,}5$ es la pendiente del \cref{eje:2-recta}.
\end{solucion}

\begin{solucion}{eje:2-recta}
\datos $x=(0,1,2)$, $y=(0,1,3)$.
\planteamiento ecuaciones normales (\cref{teo:normales}) resueltas por Cramer; después $\hat\sigma^{2}=\SSE/n$ (\cref{teo:mv-gauss}).
\pasos
\begin{enumerate}[label=\arabic*.]
  \item Filas de $X$: $(1,0)$, $(1,1)$, $(1,2)$. $X^{\top}X=\begin{pmatrix}3&0+1+2\\0+1+2&0+1+4\end{pmatrix}=\begin{pmatrix}3&3\\3&5\end{pmatrix}$.
  \item $X^{\top}y=(0+1+3,\;0+1+6)=(4,7)$.
  \item Determinante: $3\cdot 5-3\cdot 3=6$.
  \item $b=(4\cdot 5-3\cdot 7)/6=-1/6\approx-0{,}1667$; $w=(3\cdot 7-3\cdot 4)/6=9/6=3/2$.
  \item Predicciones $-1/6+1{,}5x$: $(-1/6,\;4/3,\;17/6)$.
  \item Residuos: $(1/6,\;-1/3,\;1/6)$.
  \item $\SSE=1/36+4/36+1/36=6/36=1/6\approx 0{,}1667$.
  \item $\hat\sigma^{2}=(1/6)/3=1/18\approx 0{,}05556$.
\end{enumerate}
\resultado \textbf{$\hat\theta=(b,w)=(-1/6,\;3/2)$; $\SSE=1/6$; $\hat\sigma^{2}=1/18\approx 0{,}05556$.}
\comprobacion ortogonalidad: $\sum r_i=1/6-2/6+1/6=0$ y $\sum x_ir_i=0-1/3+2/6=0$.
\end{solucion}

\begin{solucion}{eje:2-strang1}
\datos $t=(-1,1,2)$, $b=(1,3,11)$; filas de $A$: $(1,t_i)$.
\planteamiento ecuaciones normales $A^{\top}A\hat x=A^{\top}b$ y propiedad $A^{\top}e=0$ (\cref{prop:proyeccion}).
\pasos
\begin{enumerate}[label=\arabic*.]
  \item $A^{\top}A=\begin{pmatrix}3&-1+1+2\\2&1+1+4\end{pmatrix}=\begin{pmatrix}3&2\\2&6\end{pmatrix}$; $A^{\top}b=(1+3+11,\;-1+3+22)=(15,24)$.
  \item $\det=18-4=14$; $C=(15\cdot 6-2\cdot 24)/14=42/14=3$; $D=(3\cdot 24-2\cdot 15)/14=42/14=3$.
  \item (b) Recta $b=3+3t$; alturas $p=(0,6,9)$; error $e=b-p=(1,-3,2)$.
  \item (c) $A^{\top}e=(1-3+2,\;-1-3+4)=(0,0)$.
  \item (d) Con datos $e$, $A^{\top}e=0$, luego $\hat x=(A^{\top}A)^{-1}\cdot 0=0$: la recta nula, porque $e$ ya es perpendicular al espacio de columnas.
\end{enumerate}
\resultado \textbf{$(C,D)=(3,3)$; $p=(0,6,9)$; $e=(1,-3,2)$; con datos $e$ se obtiene la recta $0$.}
\comprobacion Pitágoras: $\norm{b}^{2}=1+9+121=131$ y $\norm{p}^{2}+\norm{e}^{2}=117+14=131$.
\end{solucion}

\begin{solucion}{eje:2-strang2}
\datos $t=(-2,-1,0,1,2)$, $b=(7,4,2,1,1)$.
\planteamiento $A^{\top}A=\begin{pmatrix}n&\sum t\\\sum t&\sum t^{2}\end{pmatrix}$; con $\sum t=0$ es diagonal.
\pasos
\begin{enumerate}[label=\arabic*.]
  \item $\sum t=0$ y $\sum t^{2}=10$: $A^{\top}A=\diag(5,10)$.
  \item $\sum b=15$; $\sum tb=-14-4+0+1+2=-15$; $A^{\top}b=(15,-15)$.
  \item $C=15/5=3$; $D=-15/10=-1{,}5$.
  \item Alturas $(6;\,4{,}5;\,3;\,1{,}5;\,0)$; error $e=(1;\,-0{,}5;\,-1;\,-0{,}5;\,1)$.
  \item $\SSE=1+0{,}25+1+0{,}25+1=3{,}5$.
\end{enumerate}
\resultado \textbf{$b=3-1{,}5t$ y $\SSE=3{,}5$.} $A^{\top}A$ es diagonal porque la columna de unos y la de tiempos son ortogonales cuando $\sum t_i=0$, que es lo mismo que centrar la variable.
\comprobacion $\sum e_i=0$ y $\sum t_ie_i=-2+0{,}5+0-0{,}5+2=0$.
\end{solucion}

\begin{solucion}{eje:2-laplace}
\datos $y=(1,2,4,9,14)$; ruido de Laplace con escala $s$.
\planteamiento log-verosimilitud negativa y derivadas en $b$ y en $s$.
\pasos
\begin{enumerate}[label=\arabic*.]
  \item (a) $-\ln L(b,s)=\sum_i\bigl[\ln(2s)+\abs{y_i-b}/s\bigr]=n\ln(2s)+\frac1s\sum_i\abs{y_i-b}$.
  \item (b) Para cualquier $s$, el término en $b$ es $\frac1s\sum\abs{y_i-b}$, mínimo en la mediana (\cref{eje:1-constante}): $\hat b=4$.
  \item Derivando en $s$: $n/s-\frac{1}{s^{2}}\sum\abs{y_i-\hat b}=0$, luego $\hat s=\frac1n\sum\abs{y_i-\hat b}=20/5=4$.
  \item (c) En el modelo lineal la pérdida es $\sum_i\abs{y_i-\theta^{\top}x_i}$, que no es diferenciable donde algún residuo se anula; su «gradiente» $-\sum_i\sign(r_i)x_i$ es constante a trozos, y la condición de óptimo no es un sistema lineal. Se resuelve con programación lineal o con métodos de subgradiente.
\end{enumerate}
\resultado \textbf{$-\ln L=n\ln(2s)+\frac1s\sum\abs{y_i-b}$; $\hat b=4$ y $\hat s=4$.}
\comprobacion la derivada en $s$ pasa de negativa a positiva en $s=4$: $5/3-20/9<0$ y $5/5-20/25>0$.
\end{solucion}

\begin{solucion}{eje:2-traza}
\datos $X\in\R^{n\times m}$ de rango $m$; $P=X(X^{\top}X)^{-1}X^{\top}$, simétrica e idempotente (\cref{prop:proyeccion}).
\planteamiento usar $P^{\top}=P$, $P^{2}=P$ y la propiedad cíclica de la traza, $\tr(AB)=\tr(BA)$.
\pasos
\begin{enumerate}[label=\arabic*.]
  \item (a) $(I-P)^{\top}=I-P$; $(I-P)^{2}=I-2P+P^{2}=I-P$; $PX=X(X^{\top}X)^{-1}X^{\top}X=X$, luego $(I-P)X=0$.
  \item (b) $\tr P=\tr\bigl(X(X^{\top}X)^{-1}X^{\top}\bigr)=\tr\bigl((X^{\top}X)^{-1}X^{\top}X\bigr)=\tr I_m=m$.
  \item (c) Si $Pv=\mu v$ con $v\ne 0$, entonces $\mu v=Pv=P^{2}v=\mu^{2}v$, luego $\mu^{2}=\mu$ y $\mu\in\{0,1\}$.
  \item (d) Con columna de unos y una variable, $P_{ii}=\frac1n+\frac{(x_i-\bar x)^{2}}{\sum_j(x_j-\bar x)^{2}}$. Para $x=(0,1,2,3)$: $\bar x=1{,}5$, $\sum(x_j-\bar x)^{2}=2{,}25+0{,}25+0{,}25+2{,}25=5$, y $P_{ii}=0{,}25+(2{,}25;\,0{,}25;\,0{,}25;\,2{,}25)/5=(0{,}7;\,0{,}3;\,0{,}3;\,0{,}7)$.
\end{enumerate}
\resultado \textbf{$I-P$ es un proyector ortogonal sobre el complemento de $\mathcal{C}(X)$; $\tr P=m$; los autovalores son 0 y~1.} $\blacksquare$
\comprobacion en (d), $0{,}7+0{,}3+0{,}3+0{,}7=2=m$. Por (b) y (c), $P$ tiene exactamente $m$ autovalores iguales a~1.
\end{solucion}

% ---------------------------------------------------------------------
\section{Tema 3}

\begin{solucion}{eje:3-lineal}
\datos etiquetas $y\in\{0,1\}$; modelo $\hat y=w^{\top}x+b$ ajustado por mínimos cuadrados.
\planteamiento comparar la salida de una recta con una probabilidad, y el supuesto probabilístico de cada pérdida (\cref{tab:mv}).
\pasos
\begin{enumerate}[label=\arabic*.]
  \item \textbf{Rango.} $w^{\top}x+b$ recorre todo $\R$: para entradas alejadas da valores negativos o mayores que~1, que no son probabilidades.
  \item \textbf{Pérdida.} Un ejemplo de clase~1 con $\hat y=3$ está clasificado con mucha seguridad, pero el error cuadrático $(3-1)^{2}=4$ lo penaliza; estos puntos arrastran la recta y desplazan la frontera $\hat y=1/2$.
  \item \textbf{Supuesto.} Mínimos cuadrados es máxima verosimilitud con ruido gaussiano (\cref{teo:mv-gauss}); una etiqueta binaria sigue una Bernoulli, cuya verosimilitud es la entropía cruzada binaria (\cref{prop:bce}).
\end{enumerate}
\resultado \textbf{conviene la regresión logística: salida en $(0,1)$ y pérdida de Bernoulli.}
\comprobacion $\sigma(z)\in(0,1)$ para todo $z$, y la entropía cruzada de un ejemplo bien clasificado tiende a~0 en lugar de crecer.
\end{solucion}

\begin{solucion}{eje:3-valores}
\datos $\sigma(z)=1/(1+e^{-z})$.
\planteamiento \cref{prop:sigmoide} y $e^{\pm\ln 3}=3^{\pm 1}$.
\pasos
\begin{enumerate}[label=\arabic*.]
  \item $\sigma(\ln 3)=1/(1+1/3)=3/4=0{,}75$.
  \item $\sigma(-\ln 3)=1-\sigma(\ln 3)=1/4=0{,}25$.
  \item $\sigma'(0)=\sigma(0)(1-\sigma(0))=\frac12\cdot\frac12=\frac14$.
  \item $\sigma^{-1}(0{,}8)=\ln(0{,}8/0{,}2)=\ln 4\approx 1{,}386$.
\end{enumerate}
\resultado \textbf{$3/4$, $1/4$, $1/4$ y $\ln 4\approx 1{,}386$.}
\comprobacion $\sigma(\ln 4)=1/(1+1/4)=0{,}8$.
\end{solucion}

\begin{solucion}{eje:3-gradiente}
\datos $w=(1,-2)$, $b=0{,}5$, $x=(2,1)$, $y=0$.
\planteamiento $z=w^{\top}x+b$, $p=\sigma(z)$, $\BCE$ y $\partial\BCE/\partial z=p-y$ (\cref{teo:grad-logistica}).
\pasos
\begin{enumerate}[label=\arabic*.]
  \item $z=1\cdot 2+(-2)\cdot 1+0{,}5=0{,}5$.
  \item $e^{-0{,}5}=0{,}6065$; $p=1/1{,}6065=0{,}6225$.
  \item Como $y=0$: $\ell=-\ln(1-p)=-\ln 0{,}3775=0{,}9741$.
  \item $p-y=0{,}6225$.
  \item $\nabla_w\ell=(p-y)x=0{,}6225\cdot(2,1)=(1{,}245;\,0{,}6225)$; $\nabla_b\ell=0{,}6225$.
\end{enumerate}
\resultado \textbf{$z=0{,}5$; $p\approx 0{,}6225$; $\ell\approx 0{,}974$; $\nabla_w\ell\approx(1{,}245;\,0{,}6225)$; $\nabla_b\ell\approx 0{,}6225$.}
\comprobacion el gradiente es positivo, así que un paso de descenso baja $z$ y $p$, como corresponde a un ejemplo de clase~0. Además, $\ell>\ln 2=0{,}6931$ porque $p>1/2$ apunta a la clase equivocada.
\end{solucion}

\begin{solucion}{eje:3-frontera}
\datos $b=-3$, $w=(1,2)$; puntos $A=(0,0)$, $B=(1,1)$, $C=(3,2)$.
\planteamiento la frontera es $w^{\top}x+b=0$ (\cref{prop:frontera}); la distancia de un hiperplano $w^{\top}x+b=0$ al origen es $\abs{b}/\norm{w}$.
\pasos
\begin{enumerate}[label=\arabic*.]
  \item (a) Frontera: $x_1+2x_2=3$. Distancia al origen: $3/\sqrt{1+4}=3/\sqrt5\approx 1{,}342$.
  \item (b) $z_A=-3$, $p_A=\sigma(-3)=0{,}04743$: clase~0. $z_B=1+2-3=0$, $p_B=0{,}5$: sobre la frontera (clase~1 con el convenio $p\ge 1/2$). $z_C=3+4-3=4$, $p_C=\sigma(4)=0{,}9820$: clase~1.
  \item (c) Con $(2b,2w)$ la frontera $2(x_1+2x_2-3)=0$ es la misma, pero los logits se duplican: $p_A=\sigma(-6)=0{,}002473$, $p_B=0{,}5$, $p_C=\sigma(8)=0{,}9997$. El modelo es más «seguro» sin cambiar sus decisiones.
\end{enumerate}
\resultado \textbf{frontera $x_1+2x_2=3$ a distancia $1{,}342$; $p_A=0{,}0474$, $p_B=0{,}5$, $p_C=0{,}982$; al escalar, la frontera no cambia y las probabilidades se alejan de $1/2$.}
\comprobacion el punto $(3/5,6/5)=\frac{3}{5}(1,2)$ está en la frontera ($3/5+12/5=3$) y su norma es $\frac35\sqrt5=1{,}342$. El apartado (c) anticipa la \cref{prop:separable}.
\end{solucion}

\begin{solucion}{eje:3-newton}
\datos $J(\theta)=\frac1n\sum_i\BCE(\sigma_i,y_i)$; $x_1=(1,1)$, $y_1=1$; $x_2=(1,-2)$, $y_2=0$; $\theta_0=0$.
\planteamiento gradiente y hessiano del \cref{teo:grad-logistica}; paso $\theta_1=\theta_0-H^{-1}\nabla J$.
\pasos
\begin{enumerate}[label=\arabic*.]
  \item (a) $\nabla J=\frac1n\sum(\sigma_i-y_i)x_i$ y $\nabla_\theta\sigma_i=\sigma_i(1-\sigma_i)x_i$, luego $H=\frac1n\sum\sigma_i(1-\sigma_i)x_ix_i^{\top}$, y $v^{\top}Hv=\frac1n\sum\sigma_i(1-\sigma_i)(x_i^{\top}v)^{2}\ge 0$.
  \item (b) En $\theta_0=0$: $\sigma_1=\sigma_2=0{,}5$ y $\sigma_i(1-\sigma_i)=0{,}25$.
  \item $\nabla J=\frac12\bigl[(0{,}5-1)(1,1)+(0{,}5-0)(1,-2)\bigr]=\frac12(0;\,-1{,}5)=(0;\,-0{,}75)$.
  \item $x_1x_1^{\top}+x_2x_2^{\top}=\begin{pmatrix}1&1\\1&1\end{pmatrix}+\begin{pmatrix}1&-2\\-2&4\end{pmatrix}=\begin{pmatrix}2&-1\\-1&5\end{pmatrix}$.
  \item $H=\frac12\cdot 0{,}25\cdot\begin{pmatrix}2&-1\\-1&5\end{pmatrix}=\begin{pmatrix}0{,}25&-0{,}125\\-0{,}125&0{,}625\end{pmatrix}$.
  \item $\det H=0{,}15625-0{,}015625=0{,}140625$.
  \item $H^{-1}=\frac{1}{0{,}140625}\begin{pmatrix}0{,}625&0{,}125\\0{,}125&0{,}25\end{pmatrix}=\begin{pmatrix}4{,}444&0{,}8889\\0{,}8889&1{,}778\end{pmatrix}$.
  \item $H^{-1}\nabla J=(0{,}8889\cdot(-0{,}75);\;1{,}778\cdot(-0{,}75))=(-0{,}6667;\,-1{,}333)$; $\theta_1=(2/3,\;4/3)$.
  \item (c) $J(\theta_0)=\ln 2=0{,}6931$. En $\theta_1$, $z_1=2/3+4/3=2$ y $z_2=2/3-8/3=-2$; ambos ejemplos tienen pérdida $\ln(1+e^{-2})=0{,}1269$, luego $J(\theta_1)=0{,}1269$.
\end{enumerate}
\resultado \textbf{$\nabla J=(0;\,-0{,}75)$; $H=\begin{pmatrix}0{,}25&-0{,}125\\-0{,}125&0{,}625\end{pmatrix}$; $\theta_1=(2/3,\,4/3)$; $J$ baja de $0{,}6931$ a $0{,}1269$.} $\blacksquare$
\comprobacion $H(H^{-1}\nabla J)=(0{,}25\cdot(-0{,}6667)+0{,}125\cdot 1{,}333;\;0{,}125\cdot 0{,}6667-0{,}625\cdot 1{,}333)=(0;\,-0{,}75)$. Los datos son separables, así que iteraciones sucesivas harán crecer $\norm{\theta}$ (\cref{prop:separable}).
\end{solucion}

\begin{solucion}{eje:3-separable}
\datos $x_1=-1$, $y_1=0$; $x_2=1$, $y_2=1$; modelo $\sigma(wx)$; $\lambda=0{,}1$.
\planteamiento \cref{prop:sigmoide}(iv) para la pérdida; derivadas primera y segunda para la versión regularizada.
\pasos
\begin{enumerate}[label=\arabic*.]
  \item (a) Ejemplo~1: $-\ln(1-\sigma(-w))=-\ln\sigma(w)=\ln(1+e^{-w})$. Ejemplo~2: $-\ln\sigma(w)=\ln(1+e^{-w})$. La media es $J(w)=\ln(1+e^{-w})$.
  \item $J'(w)=-\sigma(-w)<0$: $J$ es estrictamente decreciente, con $J(w)\to 0$ cuando $w\to\infty$ y $J(w)>0$ para todo $w$. No hay mínimo.
  \item (b) $G(w)=J(w)+\frac{\lambda}{2}w^{2}$; $G'(w)=-\sigma(-w)+\lambda w$; $G''(w)=\sigma(w)\sigma(-w)+\lambda>0$. $G$ es estrictamente convexa, $G'(0)=-1/2<0$ y $G'(w)\to\infty$: hay un único cero, que cumple $\lambda w^{*}=\sigma(-w^{*})$.
  \item (c) Newton sobre $g(w)=\lambda w-\sigma(-w)$, con $g'=\lambda+\sigma(w)\sigma(-w)$, desde $w=2$: $g(2)=0{,}2-0{,}1192=0{,}0808$, $g'(2)=0{,}2050$, $w=2-0{,}394=1{,}606$; después $g(1{,}606)=-0{,}0066$, $g'=0{,}2392$, $w=1{,}633$; un paso más da $w^{*}\approx 1{,}634$.
  \item $P(y_2=1)=\sigma(1{,}634)=0{,}8367$.
\end{enumerate}
\resultado \textbf{sin regularizar no hay mínimo; con $\lambda=0{,}1$, $w^{*}\approx 1{,}634$ y $P(y_2=1)\approx 0{,}837$.} $\blacksquare$
\comprobacion $\lambda w^{*}=0{,}1634$ y $\sigma(-1{,}634)=1-0{,}8367=0{,}1633$. La regularización impide que la probabilidad llegue a~1.
\end{solucion}

% ---------------------------------------------------------------------
\section{Tema 4}

\begin{solucion}{eje:4-llm}
\datos la capa de salida de un LLM (\cref{fig:cabeza}).
\planteamiento identificar cada pieza de $z=W^{\top}x$, $p=\softmax(z)$, $\ell=-\ln p_y$ (\cref{def:softmax-reg}).
\pasos
\begin{enumerate}[label=\arabic*.]
  \item $x$ es el vector oculto $h\in\R^{d_{\text{model}}}$ de la última posición, producido por el último bloque Transformer.
  \item $K=\abs{\V}$: cada clase es un token del vocabulario.
  \item $W=W_U\in\R^{d_{\text{model}}\times\abs{\V}}$, la matriz de \emph{unembedding}; cada columna es la plantilla de un token.
  \item La pérdida es $-\ln p$ del token real, la entropía cruzada de la \cref{prop:ce}.
  \item Diferencia: en el clasificador de dígitos $x$ son los píxeles, fijos; en el LLM, $h$ lo calcula la red y se aprende a la vez que $W_U$.
\end{enumerate}
\resultado \textbf{es una regresión softmax con $\abs{\V}$ clases cuya entrada $h$ aprende el Transformer.}
\comprobacion en Llama~3 8B, $W_U$ tiene $4096\cdot 128\,256=525\,336\,576$ parámetros (\cref{eje:4-llama}).
\end{solucion}

\begin{solucion}{eje:4-logits}
\datos $z=(1,2,3)$, clase correcta $y=1$.
\planteamiento softmax, $\ell=-\ln p_y=-z_y+\LSE(z)$ y $\nabla_z\ell=p-e_y$ (\cref{teo:grad-softmax}).
\pasos
\begin{enumerate}[label=\arabic*.]
  \item $e^{z}=(2{,}718;\,7{,}389;\,20{,}09)$, de suma $30{,}19$.
  \item $p=(0{,}09003;\,0{,}2447;\,0{,}6652)$.
  \item $\ell=-\ln 0{,}09003=2{,}408$.
  \item $\nabla_z\ell=p-(1,0,0)=(-0{,}9100;\,0{,}2447;\,0{,}6652)$.
  \item Con $y=3$: $\ell=-\ln 0{,}6652=0{,}4076$.
  \item Diferencia: $\ell_1-\ell_3=(-z_1+\LSE)-(-z_3+\LSE)=z_3-z_1=2$, porque el término $\LSE(z)$ es común.
\end{enumerate}
\resultado \textbf{$p\approx(0{,}0900;\,0{,}2447;\,0{,}6652)$; $\ell\approx 2{,}408$; $\nabla_z\ell\approx(-0{,}9100;\,0{,}2447;\,0{,}6652)$; con la tercera clase, $\ell\approx 0{,}408$.}
\comprobacion $p$ suma $1{,}000$ y el gradiente suma $-0{,}0001\approx 0$ por redondeo; $2{,}408-0{,}408=2{,}000$. Los valores de $p$ son los del \cref{ej:softmax} permutados, porque $(1,2,3)$ es $(2,1,0)$ invertido y desplazado en una unidad.
\end{solucion}

\begin{solucion}{eje:4-k2}
\datos $z=(1{,}3;\,0{,}4)$.
\planteamiento calcular ambos miembros; después dividir por $e^{z_1}$.
\pasos
\begin{enumerate}[label=\arabic*.]
  \item $e^{1{,}3}=3{,}669$ y $e^{0{,}4}=1{,}492$, de suma $5{,}161$; $\softmax_1=3{,}669/5{,}161=0{,}7109$.
  \item $z_1-z_2=0{,}9$; $e^{-0{,}9}=0{,}4066$; $\sigma(0{,}9)=1/1{,}4066=0{,}7109$.
  \item En general: $p_1=\dfrac{e^{z_1}}{e^{z_1}+e^{z_2}}=\dfrac{1}{1+e^{z_2-z_1}}=\sigma(z_1-z_2)$.
\end{enumerate}
\resultado \textbf{ambos valen $0{,}7109$, y $p_1=\sigma(z_1-z_2)$ para todo $z\in\R^{2}$.} $\blacksquare$
\comprobacion si $z_1=z_2$, $p_1=\sigma(0)=1/2$, como corresponde a dos clases con el mismo logit.
\end{solucion}

\begin{solucion}{eje:4-llama}
\datos $d_{\text{model}}=4096$, $\abs{\V}=128\,256$; 2 bytes por número.
\planteamiento la capa es una matriz $d_{\text{model}}\times\abs{\V}$; la pérdida inicial es $\ln K$ (\cref{prop:inicial}).
\pasos
\begin{enumerate}[label=\arabic*.]
  \item Parámetros: $4096\cdot 128\,256=525\,336\,576$.
  \item Memoria: $525\,336\,576\cdot 2=1\,050\,673\,152$ bytes $\approx 1{,}051$~GB ($\approx 0{,}9785$~GiB, dividiendo por $2^{30}$).
  \item Pérdida uniforme: $\ln 128\,256=\ln 1{,}28256+5\ln 10=0{,}2489+11{,}513=11{,}76$ nats.
\end{enumerate}
\resultado \textbf{$525\,336\,576$ parámetros; $\approx 1{,}05$~GB en bf16; pérdida inicial $\approx 11{,}76$ nats.}
\comprobacion orden de magnitud: $4\cdot 10^{3}\cdot 1{,}3\cdot 10^{5}\approx 5{,}2\cdot 10^{8}$. Para Llama~2 7B, $4096\cdot 32\,000=131\,072\,000$, cuatro veces menos, porque su vocabulario es unas cuatro veces menor.
\end{solucion}

\begin{solucion}{eje:4-gradiente}
\datos $p=\softmax(z)$, $\ell=-\ln p_y$.
\planteamiento escribir $\ell=-z_y+\LSE(z)$ y derivar cada término.
\pasos
\begin{enumerate}[label=\arabic*.]
  \item $\ell=-\ln\bigl(e^{z_y}/\sum_je^{z_j}\bigr)=-z_y+\ln\sum_je^{z_j}$.
  \item $\partial(-z_y)/\partial z_k=-\ind[k=y]$.
  \item $\frac{\partial}{\partial z_k}\ln\sum_je^{z_j}=\frac{1}{\sum_je^{z_j}}\cdot e^{z_k}=p_k$.
  \item $\partial\ell/\partial z_k=p_k-\ind[k=y]$; en forma vectorial, $\nabla_z\ell=p-e_y$.
  \item $\sum_k(p_k-\ind[k=y])=1-1=0$.
  \item Invariancia: $\ell(z+c\mathbf{1})=\ell(z)$ para todo $c$. Derivando en $c=0$, $\nabla_z\ell\cdot\mathbf{1}=0$, que es la suma nula.
\end{enumerate}
\resultado \textbf{$\nabla_z\ell=p-e_y$, con componentes de suma nula.} $\blacksquare$
\comprobacion en el \cref{eje:4-logits}, $-0{,}9100+0{,}2447+0{,}6652=-0{,}0001\approx 0$.
\end{solucion}

\begin{solucion}{eje:4-hessiano}
\datos $\ell(z)=-z_y+\LSE(z)$, con $\nabla_z\ell=p-e_y$.
\planteamiento derivar $p_k$ respecto a $z_j$ por la regla del cociente; calcular la covarianza de un vector \emph{one-hot}.
\pasos
\begin{enumerate}[label=\arabic*.]
  \item (a) Con $S=\sum_je^{z_j}$: $\partial p_k/\partial z_j=\bigl(\ind[j=k]e^{z_k}S-e^{z_k}e^{z_j}\bigr)/S^{2}=\ind[j=k]p_k-p_kp_j$. Como $e_y$ es constante, $\nabla^{2}_z\ell=\diag(p)-pp^{\top}$.
  \item (b) Sea $E$ con $P(E=e_k)=p_k$. $\E[E]=p$; $\E[EE^{\top}]=\sum_kp_ke_ke_k^{\top}=\diag(p)$; $\Cov(E)=\diag(p)-pp^{\top}$.
  \item (c) $z=(\ln 2,0,0)$: $e^{z}=(2,1,1)$, $S=4$, $p=(0{,}5;\,0{,}25;\,0{,}25)$.
  \item $pp^{\top}$ tiene entradas $0{,}25$, $0{,}125$ y $0{,}0625$; restándolo de $\diag(p)$:
  \[
  \nabla_z^{2}\ell=\begin{pmatrix}0{,}25&-0{,}125&-0{,}125\\-0{,}125&0{,}1875&-0{,}0625\\-0{,}125&-0{,}0625&0{,}1875\end{pmatrix}=\frac{1}{16}\begin{pmatrix}4&-2&-2\\-2&3&-1\\-2&-1&3\end{pmatrix}.
  \]
\end{enumerate}
\resultado \textbf{$\nabla_z^{2}\ell=\diag(p)-pp^{\top}=\Cov(E)$; en el punto dado, la matriz anterior.} $\blacksquare$
\comprobacion cada fila suma cero: $4-2-2=0$, $-2+3-1=0$, $-2-1+3=0$, de acuerdo con la \cref{prop:hess-softmax}. La diagonal es $p_k(1-p_k)$: $0{,}5\cdot 0{,}5=0{,}25$ y $0{,}25\cdot 0{,}75=0{,}1875$.
\end{solucion}

% ---------------------------------------------------------------------
\section{Tema 5}

\begin{solucion}{eje:5-vf}
\planteamiento contrastar cada afirmación con la \cref{prop:encajadas}, el \cref{teo:sesgo-varianza} y la \cref{prop:optimismo}.
\pasos
\begin{enumerate}[label=(\alph*)]
  \item \textbf{Verdadera.} Las clases de polinomios están encajadas (\cref{prop:encajadas}); en el \cref{tab:polinomios}, el ECM de entrenamiento baja de $0{,}1183$ a $0{,}0024$ y a $0{,}0000$.
  \item \textbf{Falsa.} Más regularización contrae los parámetros hacia cero y el modelo cambia menos al cambiar los datos: la varianza baja y el sesgo sube.
  \item \textbf{Falsa.} $\sigma^{2}$ es la varianza de $y$ dada $x$, propiedad de la distribución de los datos y no del modelo; más datos reducen la varianza del modelo, no el ruido.
  \item \textbf{Falsa.} El polinomio de grado~9 del \cref{tab:polinomios} tiene error de entrenamiento nulo y error de validación $34{,}17$. Además, el riesgo nunca es menor que $\sigma^{2}$.
\end{enumerate}
\resultado \textbf{V, F, F, F.}
\end{solucion}

\begin{solucion}{eje:5-diagnostico}
\datos A: $0{,}60$ en entrenamiento y $0{,}58$ en validación. B: $0{,}99$ y $0{,}70$.
\planteamiento mirar el nivel del error de entrenamiento y la distancia entre entrenamiento y validación (\cref{tab:diagnostico}).
\pasos
\begin{enumerate}[label=\arabic*.]
  \item A: error de entrenamiento alto ($40\,\%$) y diferencia pequeña (2 puntos). No aprende ni los datos que ve: \textbf{sesgo alto}. Remedio: modelo más rico, más características, menos regularización o entrenar más.
  \item B: error de entrenamiento casi nulo y diferencia de 29 puntos. Reproduce el ruido de la muestra: \textbf{varianza alta}. Remedio: más datos, más regularización, modelo más simple o parada temprana.
\end{enumerate}
\resultado \textbf{A infraajusta; B sobreajusta.}
\comprobacion B se comporta como el polinomio de grado~9 del \cref{tab:polinomios}, y A como el de grado~1.
\end{solucion}

\begin{solucion}{eje:5-mano}
\datos predicciones $1{,}8$; $2{,}4$; $2{,}1$, equiprobables; $f(x_0)=2{,}5$; $\sigma^{2}=0{,}09$.
\planteamiento \cref{teo:sesgo-varianza} con la esperanza sobre $S$ calculada como media de tres valores.
\pasos
\begin{enumerate}[label=\arabic*.]
  \item $\bar f(x_0)=(1{,}8+2{,}4+2{,}1)/3=6{,}3/3=2{,}1$.
  \item Sesgo $=2{,}1-2{,}5=-0{,}4$; sesgo$^{2}=0{,}16$.
  \item Varianza $=\frac13\bigl[(-0{,}3)^{2}+0{,}3^{2}+0^{2}\bigr]=0{,}18/3=0{,}06$.
  \item Error esperado $=0{,}09+0{,}16+0{,}06=0{,}31$.
\end{enumerate}
\resultado \textbf{sesgo$^{2}=0{,}16$, varianza $=0{,}06$, error cuadrático esperado $=0{,}31$.}
\comprobacion directamente, $\E_S[(f-\hat f)^{2}]=\frac13(0{,}7^{2}+0{,}1^{2}+0{,}4^{2})=\frac13(0{,}49+0{,}01+0{,}16)=0{,}22=0{,}16+0{,}06$; sumando el ruido, $0{,}31$.
\end{solucion}

\begin{solucion}{eje:5-interpolar}
\datos puntos $(0,1)$, $(1,0)$, $(2,3)$; polinomio $\theta_0+\theta_1x+\theta_2x^{2}$.
\planteamiento resolver $V\theta=y$ con la matriz de Vandermonde (\cref{prop:interpolacion}).
\pasos
\begin{enumerate}[label=\arabic*.]
  \item (a) $x=0$: $\theta_0=1$. $x=1$: $1+\theta_1+\theta_2=0$, es decir, $\theta_1+\theta_2=-1$. $x=2$: $1+2\theta_1+4\theta_2=3$, es decir, $\theta_1+2\theta_2=1$.
  \item Restando: $\theta_2=2$ y $\theta_1=-3$. El polinomio es $1-3x+2x^{2}$.
  \item $\det V=(x_2-x_1)(x_3-x_1)(x_3-x_2)=(1-0)(2-0)(2-1)=2$.
  \item (b) El error de entrenamiento es nulo siempre que los $x_i$ sean distintos y el grado sea al menos $n-1$: con grado $n-1$ hay solución única; con grado mayor, infinitas.
  \item (c) Si $x_i=x_j$ con $y_i\ne y_j$, ninguna función puede pasar por ambos puntos: el error de entrenamiento es positivo para cualquier grado, y $V$ tiene dos filas iguales.
\end{enumerate}
\resultado \textbf{$1-3x+2x^{2}$, con $\det V=2$; interpolación exacta si los $x_i$ son distintos y $p\ge n-1$.}
\comprobacion $1-3+2=0$ y $1-6+8=3$.
\end{solucion}

\begin{solucion}{eje:5-encoger}
\datos $x_1,\dots,x_n$ independientes con media $\mu$ y varianza $\sigma^{2}$; $\hat\mu=a\bar x$.
\planteamiento $\ECM=\text{sesgo}^{2}+\text{varianza}$ (\cref{teo:sesgo-varianza} sin ruido); minimizar en $a$.
\pasos
\begin{enumerate}[label=\arabic*.]
  \item $\E[\bar x]=\mu$, luego $\E[\hat\mu]=a\mu$ y el sesgo es $(a-1)\mu$.
  \item Por independencia, $\Var(\bar x)=\sigma^{2}/n$, luego $\Var(\hat\mu)=a^{2}\sigma^{2}/n$.
  \item (a) $\ECM=\E[(\hat\mu-a\mu)^{2}]+(a\mu-\mu)^{2}=(a-1)^{2}\mu^{2}+a^{2}\sigma^{2}/n$, porque el término cruzado tiene esperanza nula.
  \item (b) $\frac{d\,\ECM}{da}=2(a-1)\mu^{2}+2a\sigma^{2}/n=0$, luego $a^{*}=\dfrac{\mu^{2}}{\mu^{2}+\sigma^{2}/n}$. La segunda derivada, $2\mu^{2}+2\sigma^{2}/n$, es positiva.
  \item (c) $\sigma^{2}/n=1$; $a^{*}=4/5=0{,}8$; $\ECM(0{,}8)=0{,}04\cdot 4+0{,}64\cdot 1=0{,}80$, frente a $\ECM(1)=1$.
\end{enumerate}
\resultado \textbf{$a^{*}=\mu^{2}/(\mu^{2}+\sigma^{2}/n)=0{,}8$ y $\ECM=0{,}80$, frente a $1{,}00$ de la media muestral.} $\blacksquare$
\comprobacion coincide con el \cref{tab:encoger}. Si $n\to\infty$, $a^{*}\to 1$: con muchos datos no conviene encoger. El $a^{*}$ depende de $\mu$, desconocido; en la práctica el grado de contracción se elige por validación.
\end{solucion}

\begin{solucion}{eje:5-varianza}
\datos $y=X\theta+\varepsilon$, $\Cov(\varepsilon)=\sigma^{2}I$; $\hat y=Py$.
\planteamiento $\Cov(Ay)=A\Cov(y)A^{\top}$ para una matriz constante $A$; \cref{eje:2-traza}.
\pasos
\begin{enumerate}[label=\arabic*.]
  \item (a) $\Cov(\hat y)=P\,(\sigma^{2}I)\,P^{\top}=\sigma^{2}PP^{\top}=\sigma^{2}P^{2}=\sigma^{2}P$.
  \item (b) $\Var(\hat y_i)=\sigma^{2}P_{ii}$, luego $\frac1n\sum_i\Var(\hat y_i)=\frac{\sigma^{2}}{n}\tr P=\sigma^{2}m/n$.
  \item (c) La varianza media sobre los puntos de entrenamiento crece linealmente con el número de parámetros $m$: cada parámetro añade $\sigma^{2}/n$. Con $m=n$ (interpolación), alcanza $\sigma^{2}$.
  \item (d) $\Var(\hat y_i)=0{,}05\cdot(0{,}7;\,0{,}3;\,0{,}3;\,0{,}7)=(0{,}035;\,0{,}015;\,0{,}015;\,0{,}035)$, de media $0{,}025$.
\end{enumerate}
\resultado \textbf{$\Cov(\hat y)=\sigma^{2}P$ y varianza media $\sigma^{2}m/n$; en el ejemplo, $0{,}025$.} $\blacksquare$
\comprobacion $\sigma^{2}m/n=0{,}05\cdot 2/4=0{,}025$. Los puntos extremos ($x=0$ y $x=3$) tienen más varianza que los centrales, porque están más lejos de $\bar x$.
\end{solucion}

% ---------------------------------------------------------------------
\section{Tema 6}

\begin{solucion}{eje:6-limites}
\datos ridge con $b$ sin penalizar; \cref{ej:ridge}: $\hat w_\lambda=4/(2+\lambda)$.
\planteamiento con datos centrados, $\hat w_\lambda=(X_c^{\top}X_c+\lambda I)^{-1}X_c^{\top}y_c$ y $\hat b=\bar y-\hat w^{\top}\bar x$.
\pasos
\begin{enumerate}[label=\arabic*.]
  \item $\lambda\to 0^{+}$: $\hat w_\lambda$ tiende a la solución de mínimos cuadrados (a la de norma mínima si $X$ no tiene rango completo).
  \item $\lambda\to\infty$: $\hat w_\lambda\to 0$ y $\hat b\to\bar y$; el modelo predice la media de las respuestas.
  \item Cuarta parte de $\hat w=2$: $4/(2+\lambda)=0{,}5$, luego $2+\lambda=8$ y $\lambda=6$.
\end{enumerate}
\resultado \textbf{$\lambda\to 0^{+}$: mínimos cuadrados; $\lambda\to\infty$: $\hat w\to 0$ y $\hat b\to\bar y$; en el ejemplo, $\lambda=6$.}
\comprobacion con $\lambda=6$, $\mathrm{df}=2/8=0{,}25$: el número efectivo de parámetros también se reduce a la cuarta parte.
\end{solucion}

\begin{solucion}{eje:6-ridge}
\datos $x=(0,1,2)$, $y=(0,1,3)$, $\lambda=2$ solo sobre $w$.
\planteamiento con datos centrados, $\hat w=\sum x_{c,i}y_{c,i}/(\sum x_{c,i}^{2}+\lambda)$ y $\hat b=\bar y-\hat w\bar x$.
\pasos
\begin{enumerate}[label=\arabic*.]
  \item $\bar x=1$; $\bar y=4/3$.
  \item $x_c=(-1,0,1)$; $y_c=(-4/3,\;-1/3,\;5/3)$.
  \item $\sum x_cy_c=4/3+0+5/3=3$; $\sum x_c^{2}=2$.
  \item $\hat w=3/(2+2)=3/4$ (sin regularizar sería $3/2$).
  \item $\hat b=4/3-3/4=7/12\approx 0{,}5833$.
  \item Comprobación matricial: $\begin{pmatrix}3&3\\3&5+2\end{pmatrix}\theta=\begin{pmatrix}4\\7\end{pmatrix}$, con determinante $21-9=12$: $b=(28-21)/12=7/12$, $w=(21-12)/12=3/4$.
\end{enumerate}
\resultado \textbf{$\hat w=0{,}75$ y $\hat b=7/12\approx 0{,}5833$}: la pendiente se reduce a la mitad.
\comprobacion la recta pasa por $(\bar x,\bar y)$: $7/12+3/4=16/12=4/3$, como exige la observación sobre el término independiente no penalizado.
\end{solucion}

\begin{solucion}{eje:6-cv}
\datos $n=1079$, $k=5$, cuatro valores de $\lambda$ y un entrenamiento final.
\planteamiento repartir $n$ en $k$ partes lo más iguales posible (\cref{def:kfold}).
\pasos
\begin{enumerate}[label=\arabic*.]
  \item $1079=5\cdot 215+4$: cuatro partes de $216$ y una de $215$.
  \item Se entrena con $1079-216=863$ ejemplos, o con $1079-215=864$ cuando se valida con la parte pequeña.
  \item Entrenamientos de la validación cruzada: $5\cdot 4=20$; con el final, $21$.
\end{enumerate}
\resultado \textbf{partes de 216, 216, 216, 216 y 215; 863 u 864 ejemplos por ronda; 21 entrenamientos.}
\comprobacion $4\cdot 216+215=864+215=1079$: cada ejemplo se valida exactamente una vez por valor de $\lambda$.
\end{solucion}

\begin{solucion}{eje:6-lasso}
\datos $\hat\beta=(3;\,-0{,}5;\,1{,}2)$, $X^{\top}X=I$, $\lambda=1$.
\planteamiento \cref{prop:umbral}: lasso $\sign(\hat\beta_j)(\abs{\hat\beta_j}-\lambda)_{+}$; ridge $\hat\beta_j/(1+\lambda)$.
\pasos
\begin{enumerate}[label=\arabic*.]
  \item Lasso: $3-1=2$; $\abs{-0{,}5}=0{,}5\le 1$, luego $0$; $1{,}2-1=0{,}2$. Solución $(2;\,0;\,0{,}2)$.
  \item Ridge: $\hat\beta/2=(1{,}5;\,-0{,}25;\,0{,}6)$.
  \item El mayor coeficiente, 3, pierde una unidad con el lasso ($-33\,\%$) y la mitad con ridge ($-50\,\%$).
\end{enumerate}
\resultado \textbf{lasso $(2;\,0;\,0{,}2)$, que anula el segundo coeficiente; ridge $(1{,}5;\,-0{,}25;\,0{,}6)$, que no anula ninguno.}
\comprobacion para el segundo coeficiente, $g(t)=\frac12(t+0{,}5)^{2}+\abs{t}$ vale $0{,}125$ en $t=0$, $0{,}1513$ en $t=-0{,}05$ y $0{,}2013$ en $t=0{,}05$: el mínimo está en~0.
\end{solucion}

\begin{solucion}{eje:6-map}
\datos $y\mid\theta\sim\N(X\theta,\sigma^{2}I)$; $\theta\sim\N(0,\tau^{2}I)$.
\planteamiento por Bayes, $p(\theta\mid y)\propto p(y\mid\theta)\,p(\theta)$; se minimiza $-\ln p(\theta\mid y)$.
\pasos
\begin{enumerate}[label=\arabic*.]
  \item El denominador $p(y)$ no depende de $\theta$ y no afecta al maximizador.
  \item $-\ln p(y\mid\theta)=\frac{1}{2\sigma^{2}}\norm{y-X\theta}^{2}+\text{cte}$ (\cref{teo:mv-gauss}).
  \item $-\ln p(\theta)=\frac{1}{2\tau^{2}}\norm{\theta}^{2}+\text{cte}$.
  \item $-\ln p(\theta\mid y)=\frac{1}{2\sigma^{2}}\norm{y-X\theta}^{2}+\frac{1}{2\tau^{2}}\norm{\theta}^{2}+\text{cte}$.
  \item Multiplicando por $2\sigma^{2}>0$: $\norm{y-X\theta}^{2}+\frac{\sigma^{2}}{\tau^{2}}\norm{\theta}^{2}$, que es $J_\lambda$ con $\lambda=\sigma^{2}/\tau^{2}$.
  \item $\tau\to\infty$: la distribución \emph{a priori} se vuelve plana, $\lambda\to 0$ y se recupera la máxima verosimilitud (mínimos cuadrados).
  \item $\tau\to 0^{+}$: la creencia en $\theta=0$ es cada vez más firme, $\lambda\to\infty$ y $\hat\theta\to 0$.
\end{enumerate}
\resultado \textbf{MAP $=$ ridge con $\lambda=\sigma^{2}/\tau^{2}$.} $\blacksquare$
\comprobacion en el \cref{ej:ridge}, con $\sigma^{2}=1$ y $\tau^{2}=0{,}5$ se obtiene $\lambda=2$ y $\hat w=4/(2+2)=1$, frente a $\hat w=2$ de máxima verosimilitud.
\end{solucion}

\begin{solucion}{eje:6-ruido}
\datos $X$ fija; $E$ con entradas $e_{ij}$ independientes, $\E[e_{ij}]=0$, $\E[e_{ij}^{2}]=s^{2}$.
\planteamiento calcular las esperanzas entrada a entrada usando la independencia.
\pasos
\begin{enumerate}[label=\arabic*.]
  \item (a) $(X+E)^{\top}(X+E)=X^{\top}X+X^{\top}E+E^{\top}X+E^{\top}E$. Los términos cruzados tienen esperanza nula porque $\E[E]=0$.
  \item $(E^{\top}E)_{jk}=\sum_ie_{ij}e_{ik}$, con $\E[e_{ij}e_{ik}]=s^{2}$ si $j=k$ y $0$ si $j\ne k$. Luego $\E[E^{\top}E]=ns^{2}I$.
  \item (b) $y-(X+E)\theta=(y-X\theta)-E\theta$. Entonces $\E\norm{y-(X+E)\theta}^{2}=\norm{y-X\theta}^{2}-2(y-X\theta)^{\top}\E[E]\theta+\theta^{\top}\E[E^{\top}E]\theta=\norm{y-X\theta}^{2}+ns^{2}\norm{\theta}^{2}$.
  \item Es la función objetivo ridge con $\lambda=ns^{2}$.
  \item (c) En el \cref{ej:ridge}, $n=3$ y $\lambda=2$, luego $s^{2}=2/3\approx 0{,}6667$.
\end{enumerate}
\resultado \textbf{perturbar las entradas con ruido de varianza $s^{2}$ equivale, en media, a ridge con $\lambda=ns^{2}$; en el ejemplo, $s^{2}=2/3$.} $\blacksquare$
\comprobacion con $s^{2}=2/3$, $\E\bigl[\sum_i(x_i+e_i)^{2}\bigr]=2+3\cdot 2/3=4=\sum x_i^{2}+\lambda$, el denominador de $\hat w_\lambda=4/(2+\lambda)$.
\end{solucion}
